\section{Introduction}
\label{sec:intro}
An established quantized language-model deployment already has a numerical behavior: its prompt processing creates a particular KV state, and its continuation produces particular target vectors and tokens. Parallel execution is attractive because many inputs are known together. A prompt provides many query rows; speculative decoding proposes several continuation tokens for one target verification~\cite{leviathan2023,chen2023}. Preserving the deployment's behavior allows these optimizations to reuse its regression expectations and stateful continuations. The question is how to change the organization of this work while retaining its outputs and state.

A four-query Llama execution illustrates the difficulty. The single-query reference consumed physical KV extents of 768, 768, 768 and 1024 entries. A wider call consumed 1024 for every query. The visible attention inputs and probability prefixes agreed, and the added probabilities were exactly zero, yet the attention output differed. Replaying the captured operands with cropped or padded physical ranges reproduced the two results. The first saved target-vector discrepancy appeared well before the generated tokens diverged. Thus, matching the model, positions and causal mask does not determine the computation performed by a quantized GPU runtime.

We investigate this problem through interventions on native execution. The physical-range experiment identifies where attention must be split. Projection experiments identify the activation and reduction organization that wider execution must retain. A placement experiment shows that width can also move an operation between CPU and GPU. Producer and commit-state experiments then show why preserving one operator in isolation is insufficient: the incoming state and the ownership of its resulting vectors must remain compatible across phase boundaries. Together, these observations provide the requirements for an execution method.

The method constructs a \emph{reference-derived phase plan} from a backend arithmetic recipe and the live cache geometry. Source analysis and native operand replays establish the recipe before deployment. At each call, the planner reads occupancy, physical head, capacity and padding to assign each query its reference extent. Adjacent queries with compatible requirements form a group. Whole-graph execution places cuts around the complete computation; local execution confines them to attention and keeps projection and feed-forward work wider. Whole execution is both a usable implementation and the strong simple control for evaluating local cuts.

A separate choice exposes parallelism inside the wider projection work. On the affected CUDA Q8 path, independent groups retain the native four-column computation and run as channels of one GPU operation. A mixed Q4\_K/Q6\_K extension retains scalar-column arithmetic and exposes the columns as independent channels. These arithmetic groups are distinct from the query groups defined by physical KV extent. The executor associates requested outputs with logical rows and maintains the committed frontier on one resident target/KV owner. Prefill, continuation and verification use the same state interface with their original phase-specific numerical behavior.

Existing work provides several foundations for this construction. Batch-invariant operators control reduction choices~\cite{thinkingmachines}; arithmetic reconstruction and compiler enforcement preserve existing library behavior~\cite{taming}; deterministic engines connect kernel properties with state invariants~\cite{vosti}. Fixed-shape schedules and verified speculation provide system designs for consistent execution~\cite{corun,llm42}. Our focus is the hardware-aware organization of known-query work under a quantized deployment fixed before optimization. Its physical cache range, native arithmetic and phase state determine where work can be combined.

The evaluation first checks full target vectors and active state, then separates the effects of cut placement and arithmetic packing. The implementations cover CUDA and Vulkan, three model families and a 32B mixed-quantization path. In balanced Linux/A5000 experiments, Q8 packing contributes 0.80\%--1.64\%; the contribution of local cuts is smaller because most natural verification calls do not cross a physical boundary. Compatible speculative applications reduce resident service time by 17.21\%, 25.54\% and 11.52\% for Qwen3-8B, Qwen3-14B and Llama-8B against their fastest measured fixed target-only controls. Those application gains include speculation itself. Equal-work and live-request experiments further show why a cheaper wider target call can still produce a slower request when it induces more draft and verification work.

The paper makes two contributions. First, native interventions identify how physical attention extent, backend arithmetic and producer/state history constrain changes to query width. Second, we construct an execution method that derives query cuts from those requirements, parallelizes independent arithmetic groups and composes the resulting calls through explicit output and state ownership. Numerical validation and controlled service measurements establish the behavior and practical value of that construction.
